\documentclass[12pt, a4paper]{article}

% Paquetes requeridos
\usepackage[utf8]{inputenc}
\usepackage[T1]{fontenc}
\usepackage[spanish]{babel}
\usepackage{geometry}
\usepackage{amsmath}
\usepackage[american]{circuitikz}
\usetikzlibrary{babel} % Evita conflictos entre las flechas de circuitikz y babel-spanish

% Configuracion de la pagina
\geometry{a4paper, margin=2.5cm}

\begin{document}

\section*{Ejemplo 4.1: Modelo con Fuente Dependiente de Corriente}

\subsection*{Descripción del Circuito}
El siguiente esquema eléctrico modela un sistema de dos puertos con las siguientes características:
\begin{itemize}
    \item Entre los nodos \textbf{a} y \textbf{b} se encuentra una fuente independiente de tensión ($V_g$) en paralelo con una resistencia $R_g$.
    \item La corriente que circula por la resistencia $R_g$ se define como $i_1$.
    \item Entre los nodos \textbf{c} y \textbf{b} existe una fuente dependiente de corriente cuyo valor es $\beta i_1$.
    \item En paralelo con la fuente dependiente, también entre los nodos \textbf{c} y \textbf{b}, se tiene una rama compuesta por una resistencia $R_c$ en serie con una fuente independiente de tensión ($V_{cc}$).
    \item La tensión de salida del sistema ($V_{out}$) se establece entre los terminales abiertos correspondientes a los nodos \textbf{c} y \textbf{b}.
\end{itemize}

\begin{figure}[h!]
    \centering
    \begin{circuitikz}[american, scale=1.2, transform shape]
        
        % Raíl inferior (Nodo común 'b')
        \draw (0,0) -- (7.5,0) node[ocirc] {}; 
        \node[below] at (3.5,0) {\textbf{b} (Referencia común)};

        % --- PUERTO DE ENTRADA ---
        % Nodos de entrada
        \draw (0,0) to[V, l=$V_g$] (0,3) -- (2,3) node[circ] (A) {};
        \node[above] at (A) {a};
        
        % Resistencia Rg y corriente i1
        \draw (A) to[R, l=$R_g$, i=$i_1$] (2,0);
        
        % --- PUERTO DE SALIDA ---
        % Nodo c
        \node[circ] (C) at (4,3) {};
        \node[above] at (C) {c};
        
        % Fuente de corriente dependiente
        \draw (C) to[cI, l=$\beta i_1$] (4,0);
        
        % Rama Rc - Vc en serie
        \draw (C) -- (6,3) node[circ] {} to[R, l=$R_c$] (6,1.5) to[V, l=$V_{cc}$] (6,0);
        
        % Terminales de salida
        \draw (6,3) -- (7.5,3) node[ocirc] {};
        \draw[<->] (7.5,2.8) -- (7.5,0.2) node[midway, right] {$V_{out}$};

    \end{circuitikz}
    \caption{Esquema eléctrico del Ejemplo 4-1.}
    \label{fig:ejm-4-1}
\end{figure}

\end{document}